% !TeX root = ../main.tex
\section{Datenbeschaffung und explorative Datenanalyse}
\label{sec:daten}

\subsection{Datenquelle und Erhebung}
Der verwendete Scraper entstand 2025 im Rahmen eines anderen Projekts. Mit
Python und Selenium wurden Mitte 2025 auf
\href{https://www.swimrankings.net}{swimrankings.net} die TOP100 Listen mit
bis zu 100 Männern und 100 Frauen je Land sowie deren persönliche
Bestleistungen erfasst. Die vorliegende
Arbeit verwendet diese gespeicherten Rohdaten weiter. Eine Zeile entspricht
einem Bestleistungseintrag und verbindet Athletenmerkmale wie Jahrgang,
Geschlecht, Nation und Verein mit Disziplin, Beckenlänge, Zeit, Punkten sowie
Datum, Ort und Wettkampfname.

Das Land der ausgewählten Liste (Listenland) kann von der Nation im
Athletenprofil (Profilnation) abweichen. Beispielsweise erscheinen Athleten
mit spanischer Profilnation auch auf portugiesischen Listen. Beim
Zusammenführen aller Länderlisten können deshalb mehr als 200 Personen
derselben Profilnation im Datensatz enthalten sein.

Ein ergänzendes Skript sollte 126 vermutete Erhebungsausfälle nachladen.
Bei diesen Profilen fehlten sowohl Bestleistungen als auch Detailangaben.
Der im EDA Notebook dokumentierte Versuch scheiterte an Cloudflare
Prüfungen, obwohl das Skript mit einer manuellen Prüfung im Browser durchgeführt wurde. Es wurden dadurch keine zusätzlichen Daten gewonnen.

\subsection{Bereinigung und Datenqualität}
Entfernt wurden 1\,688 Zeilen ohne Bestleistungen von 1\,687 Athleten,
134 exakte Duplikate und der unplausible Quellenwert \texttt{55:55.55}
über 200\,m Brust. Fehlende Bestleistungen gelten dabei nicht grundsätzlich
als Erhebungsfehler. Zeiten wurden in Sekunden umgerechnet, Datentypen
vereinheitlicht und fehlende Angaben erkennbar belassen. Fehlende Orte und
Wettkampfnamen sind als \texttt{Unknown} markiert. Die kombinierte Profilangabe
wurde in Nation und Verein aufgeteilt. Listenland und Profilnation bleiben
erhalten, da sie bei 215 Athleten voneinander abweichen.

Quellenstichproben und Plausibilitätsprüfungen ergänzten die Bereinigung.
Die auffälligen ungefähren Alterswerte von 6 und 66 Jahren wurden bestätigt und
beibehalten. Der Parquet Export wurde erneut eingelesen und auf
Übereinstimmung mit der bereinigten Tabelle geprüft. Tabelle~\ref{tab:datenumfang}
zeigt den resultierenden Datenumfang. Die einzelnen Prüfungen sind in
\nolinkurl{Notebooks/01_eda.ipynb} dokumentiert.

\begin{table}[H]
  \centering
  \caption{Datenumfang vor und nach der Bereinigung. Eigene Auswertung.}
  \label{tab:datenumfang}
  \begin{tabular}{@{}lrr@{}}
    \toprule
    Datenstand & Zeilen & Athleten \\
    \midrule
    Rohdaten & 286\,996 & 14\,849 \\
    Bereinigt & 285\,173 & 13\,162 \\
    \bottomrule
  \end{tabular}
\end{table}

\subsection{Erkenntnisse und Grenzen}
Für die Auswertung nach Nation zählt jede Athletenkennung einmal und wird
der Profilnation zugeordnet. Die 13\,162 Athleten verteilen sich auf
182 Profilnationen. Die Abdeckung
reicht von einer bis zu 218 Personen pro Nation, der Median beträgt nur 28.
Männer machen 55,7\,\% und Frauen 44,3\,\% aus. Die 126 vermuteten
Erhebungsausfälle betreffen 122 Frauen und 4 Männer. Länderanalysen und
Geschlechtervergleiche können dadurch verzerrt sein. Bei den bekannten
Geburtsjahrgängen liegt der Median bei 1997, die mittleren 50\,\% liegen
zwischen 1989 und 2002.

Pro Athlet sind im Median 19 Bestleistungseinträge gespeichert, die mittleren
50\,\% umfassen 9 bis 34 Einträge. Einzelne Personen sind damit deutlich
umfangreicher dokumentiert als andere. Die Leistungen verteilen sich mit
52,0\,\% auf die Langbahn (50\,m) und mit 48,0\,\% auf die Kurzbahn (25\,m).
Neben vollständigen Strecken sind auch Zwischenzeiten mit dem Zusatz
\texttt{Lap} enthalten.

Die datierten Einträge reichen vom 17. Oktober 1964 bis zum 8. April 2025.
Davon stammen 94,7\,\% aus den Jahren ab 2005. Abbildung~\ref{fig:eda}
zeigt den lokalen Anstieg 2009 und den Einbruch 2020. Ursachen lassen sich
daraus allein nicht ableiten. Das Jahr 2025 ist nur teilweise erfasst und
daher nicht direkt mit vollständigen Jahren vergleichbar. Für die weitere
Modellierung bleiben alle Jahre erhalten.

\begin{figure}[H]
  \centering
  \includegraphics[width=\textwidth]{02_eda_ueberblick.pdf}
  \caption{Zeitliche Verteilung und Umfang der gespeicherten Bestleistungen.
  Links fehlt ein Eintrag ohne Datum. Rechts zählt jeder Athlet einmal.
  Eigene Auswertung.}
  \label{fig:eda}
\end{figure}

Bei 5\,115 Athleten (38,9\,\%) fehlen Vereinsangaben. Auch die Punkte sind
nicht vollständig: Von 53\,577 Einträgen ohne Punkte entfallen 53\,019 auf
\texttt{Lap} und 558 auf seltene Kombinationen aus Disziplin und Beckenlänge.
Diese Lücken bestehen bereits in der Quelle und häufen sich somit bei
bestimmten Leistungsarten.

Die Auswahl über TOP100 Listen und persönliche Bestleistungen bildet weder
den gesamten Schwimmsport noch vollständige Wettkampfbiografien ab. Die
Häufigkeiten beschreiben gespeicherte Bestleistungen und erlauben keine
direkten Aussagen zur Anzahl aller Wettkampfteilnahmen. Gekürzte und
wiederholte Wettkampfnamen bleiben erhalten, sind für sich allein aber keine
eindeutigen Veranstaltungskennungen. Die Datenbasis bleibt auf dem
Erhebungsstand von 2025.
